% !TEX root = ../main.tex
\documentclass[../main.tex]{subfiles}

\begin{document}

\chapter{模板演示}

这一章展示模板的层级、主要数学环境和个人评述环境。正式写作时，每一章都放在自己的 `chapN/chap.tex` 中，并可用 `subfiles` 单独编译。

\section{环境一览}

\begin{description}
  \item[definition] 共享定理编号。概念整理，不必逐句翻译原书定义。
  \item[assumption] 共享定理编号。模型条件、空间设定、正则性条件。
  \item[theorem] 共享定理编号。主要结论和收敛率。
  \item[lemma] 共享定理编号。技术引理、证明中间件。
  \item[corollary] 共享定理编号。由定理直接推出的结论。
  \item[proposition] 共享定理编号。重要性介于定理和引理之间的命题。
  \item[example] 共享定理编号。例子、反例、具体 prox 计算。
  \item[algorithm] 共享定理编号。算法步骤，适合非浮动排版。
  \item[exercise] 共享定理编号。练习、思考题，使用浅灰色块。
  \item[remark] 共享定理编号。轻量说明、约定、局部观察。
  \item[proof] 无独立编号。证明环境，建议写在对应命题色块内部。
  \item[solution] 无独立编号。解答环境，建议写在对应练习或例子色块内部。
  \item[note] 单独编号。你的补充、评述、疑问和延伸阅读。
\end{description}

\begin{note}[编号策略]
除 `note` 外，数学环境共享同一套章内编号。例如“定义 1.1”之后可以是“假设 1.2”和“定理 1.3”。这样更贴近数学书写，也便于交叉引用。`note` 单独编号，避免把个人评述混进书稿主线。
\end{note}

\section{书稿主线环境}

\begin{definition}[扩展实值函数]
设 $E$ 是有限维实内积空间。若函数 $f$ 的取值属于 $(-\infty,+\infty]$，则称 $f$ 为扩展实值函数。它的有效定义域记为
\[
  \dom f = \set{x\in E: f(x)<+\infty}.
\]
\end{definition}

\begin{assumption}[基本空间]
本文默认 $E$ 与 $V$ 都是有限维实内积空间，并配备内积 $\ip{\cdot}{\cdot}$ 与范数 $\norm{\cdot}$。
\end{assumption}

\begin{theorem}[投影的一阶最优性条件]\label{thm:projection}
设 $C\subseteq E$ 非空、闭且凸，$x\in E$。若
$p=\argmin_{y\in C}\frac12\norm{y-x}^2$，则
\[
  \ip{x-p}{y-p}\le 0,\quad \forall y\in C.
\]

\begin{proof}
对任意 $y\in C$ 和 $t\in[0,1]$，点 $p+t(y-p)$ 仍在 $C$ 中。由最优性，函数
$t\mapsto \frac12\norm{p+t(y-p)-x}^2$ 在 $t=0$ 处右导数非负，于是 $\ip{p-x}{y-p}\ge 0$，等价于结论。
\end{proof}
\end{theorem}

\begin{lemma}[强凸性给出二次下界]
若 $f$ 是 $\mu$-强凸函数，且 $x^\star$ 是它的极小点，则
\[
  f(x)-f(x^\star)\ge \frac{\mu}{2}\norm{x-x^\star}^2.
\]
\end{lemma}

\begin{corollary}[强凸极小点唯一]
若 $f$ 是 $\mu$-强凸函数且存在极小点，则该极小点唯一。
\end{corollary}

\begin{proposition}[近端点的最优性条件]
设 $g$ 是 proper、closed、convex 函数，$\gamma>0$。若
\[
  z=\prox_{\gamma g}(x),
\]
则
\[
  \frac{1}{\gamma}(x-z)\in \subdiff g(z).
\]
\end{proposition}

\begin{example}[$\ell_1$ 正则项的 prox]
当 $g(x)=\lambda\norm{x}_1$ 时，$\prox_{\gamma g}$ 是逐坐标软阈值算子：
\[
  \prox_{\gamma g}(x)_i
  =
  \operatorname{sign}(x_i)\max\{\abs{x_i}-\gamma\lambda,0\}.
\]
\end{example}

\begin{exercise}
验证当 $g=\delta_C$ 时，近端梯度法退化为投影梯度法。这里 $\delta_C$ 表示集合 $C$ 的示性函数。

\begin{solution}
由近端算子的定义，
\[
  \prox_{\gamma\delta_C}(z)=\argmin_{y\in C}\frac{1}{2}\norm{y-z}^2=P_C(z).
\]
因此令 $z=x^k-\gamma\grad f(x^k)$，近端梯度步变为
\[
  x^{k+1}=P_C\bigl(x^k-\gamma\grad f(x^k)\bigr),
\]
这正是投影梯度法的更新。
\end{solution}
\end{exercise}

\begin{algorithm}[近端梯度法]
考虑复合优化问题
\[
  \min_x F(x)\equiv f(x)+g(x),
\]
其中 $f$ 可微，$g$ 可做 prox 计算。
\begin{enumerate}[label=(\alph*)]
  \item 初始化 $x^0$，选择步长 $\gamma>0$。
  \item 对 $k=0,1,2,\ldots$，计算
  \[
    x^{k+1}=\prox_{\gamma g}\bigl(x^k-\gamma\grad f(x^k)\bigr).
  \]
  \item 根据目标值、梯度映射或迭代差停止。
\end{enumerate}
\end{algorithm}

\begin{remark}[交叉引用]
\Cref{thm:projection} 展示了 `cleveref` 的中文引用效果。正式写作时，建议给关键结论加标签，普通例子则不必过度标注。
\end{remark}

\section{个人评述环境}

\begin{note}[不是翻译，而是评述]
这里写你自己的理解、补充证明、与其他教材的对照、计算细节、疑问和后续要查的文献。这个环境不使用厚重边框，只保留左侧细线，作用是提示“这是我的话”，而不是装饰页面。
\end{note}

\begin{note}[proximal gradient 的关系图]
读到 proximal gradient 时，可以在这里补充：
\begin{enumerate}
  \item 它是在最小化 $f$ 的二次上界加上 $g$；
  \item 当 $g=\delta_C$ 时，它退化为 projected gradient；
  \item 当 $g=\lambda\norm{\cdot}_1$ 时，它给出 ISTA。
\end{enumerate}
这些内容不是原书逐句翻译，而是自己的知识组织。
\end{note}

\end{document}
